\section{El sistema completo: población, sustrato y ambiente controlado}

\subsection{De una larva a una población}

El modelo de Eriksen describe la dinámica de los compartimentos de una larva
individual. Un sistema de bioconversión real contiene una población de
$N(t)$ larvas, cuyas tasas agregadas determinan los flujos de materia y
energía a escala del contenedor.

\begin{definition}[Tasa agregada de población]\label{def31}
Sea $r(t)$ una tasa por larva (asimilación, síntesis, producción de CO$_2$)
y $N(t)$ el número de larvas vivas en el contenedor. La \emph{tasa agregada}
de la población es
\begin{equation}
R(t) = N(t)\,r(t),
\label{eq:agregada}
\end{equation}
bajo el supuesto de población homogénea: todas las larvas comparten el mismo
estado $(A, B, L)$, es decir, la misma edad, peso y condición fisiológica.
\end{definition}

\begin{remark}\label{rem:homogeneidad}
El supuesto de población homogénea es razonable en crías sincronizadas
(larvas neonatas sembradas el mismo día, práctica estándar en producción),
pero se degrada conforme avanza el ciclo: la variabilidad individual en las
tasas de crecimiento genera una distribución de pesos que se ensancha con el
tiempo. Una refinamiento posible ---que dejamos como extensión--- es
reemplazar \eqref{eq:agregada} por la integral sobre la distribución de
pesos de la población, $R(t) = \int r(w)\,n(w,t)\,dw$, a costa de convertir
el modelo en una ecuación de estructura poblacional. En lo que sigue se
adopta \eqref{eq:agregada}, y la mortalidad se admite mediante $N(t)$
decreciente si los datos lo requieren.
\end{remark}

\subsection{El contenedor como volumen de control}

\begin{definition}[Sistema de bioconversión instrumentado]\label{def32}
El \emph{sistema completo} es el volumen de control delimitado por las
paredes del contenedor de cría, cuya frontera es atravesada por dos flujos
de aire ---\emph{inlet} y \emph{outlet}--- y en t=0  por la carga inicial del lote 
(larvas neonatas y sustrato). Su vector de estado se organiza en
tres grupos:
\begin{equation}
\mathbf{x} =
\underbrace{(A,\, B,\, L)}_{\text{larva promedio}}
\;\cup\;
\underbrace{(N,\, S)}_{\text{población y sustrato}}
\;\cup\;
\underbrace{(T,\, w,\, c_{\mathrm{CO_2}},\, c_{\mathrm{CH_4}})}_{\text{ambiente del aire}},
\label{eq:estado-completo}
\end{equation}
donde $S$ es la masa de sustrato disponible, $T$ la temperatura del aire,
$w$ la humedad absoluta y $c_i$ las concentraciones de los gases medidos.
\end{definition}

\begin{remark}\label{rem:frontera}
La elección de la frontera determina qué es flujo y qué es transformación
(cf.~Fig.~\ref{fig:balance}). Con la frontera en las paredes del contenedor,
la alimentación y la ventilación son flujos $J_{\mathrm{in}}$,
$J_{\mathrm{out}}$, mientras que el consumo de sustrato, el crecimiento
larval, la respiración, la evaporación y la metanogénesis son
transformaciones internas $P$, $C$. El aire del \emph{inlet} aporta O$_2$,
vapor de agua y calor sensible a condiciones controladas
$(T_{\mathrm{in}}, w_{\mathrm{in}})$; el \emph{outlet} remueve aire
enriquecido en CO$_2$, CH$_4$ y vapor, y es precisamente en esa corriente de
salida donde se instalan las mediciones
$\mathbf{y} = (T, H, \mathrm{CO_2}, \mathrm{CH_4})$.
\end{remark}

\begin{remark}[Operación por lote]\label{rem:batch}
El sistema opera \emph{por lote}: en $t = 0$ se siembra una población de
$N_0$ larvas neonatas sobre una masa inicial $S_0$ de sustrato, se cierra el
contenedor y el ciclo transcurre sin recargas hasta que las larvas alcanzan
el estado de prepupa. No hay eventos discretos intermedios: la dinámica es
continua en todo el horizonte $[0, t_f]$, con condiciones iniciales
$\mathbf{x}(0)$ conocidas. Además de las mediciones de gas, el contenedor
está montado sobre una celda de carga que registra la masa total del lote,
$m_{\mathrm{LC}}(t) = S(t) + N(t)\bigl(A(t) + B(t) + L(t)\bigr) + m_f$,
donde $m_f$ es la tara (contenedor y estructura). Como la única vía de
pérdida de masa es la corriente de salida, $m_{\mathrm{LC}}$ debe decrecer
monótonamente a una tasa igual al flujo neto de masa por el \emph{outlet}
---una restricción de consistencia global que servirá de validación
independiente de los balances de gases.
\end{remark}

\begin{figure}[htbp]
\centering
\begin{tikzpicture}[
  font=\small,
  cont/.style={draw, very thick, rounded corners=6pt, minimum width=11.0cm,
               minimum height=4.8cm, fill=gray!5},
  flujo/.style={-{Stealth[length=3mm]}, very thick},
  comp/.style={draw, thick, rounded corners=3pt, fill=blue!8,
               minimum width=3.3cm, minimum height=1.35cm, align=center},
  sens/.style={draw, thick, fill=red!12, rounded corners=2pt,
               minimum width=1.0cm, minimum height=0.65cm, align=center},
  lbl/.style={fill=white, inner sep=1.6pt, font=\footnotesize, align=center},
  lblf/.style={above=2pt, midway, font=\footnotesize, align=center}
]

% Contenedor
\node[cont] (box) {};
\node[anchor=north west, xshift=8pt, yshift=-6pt] at (box.north west)
  {\textbf{Contenedor (volumen de control)}};

% Compartimentos internos (más arriba)
\node[comp] (sus) at ([xshift=-3.6cm,yshift=-0.45cm]box.center)
  {Sustrato\\[2pt] $S$};
\node[comp, fill=orange!12] (lar) at ([xshift=0cm,yshift=-0.45cm]box.center)
  {Población larval\\[2pt] $N$\\ $(A,B,L)$};
\node[comp, fill=green!8] (air) at ([xshift=3.6cm,yshift=-0.45cm]box.center)
  {Aire interno\\[2pt] $T,\, w$\\ $c_{\mathrm{CO_2}},\, c_{\mathrm{CH_4}}$};

% Intercambios internos
\draw[{Stealth[length=2.2mm]}-{Stealth[length=2.2mm]}, thick, orange!85!black]
  (sus) -- node[lbl, midway]{consumo} (lar);
\draw[{Stealth[length=2.2mm]}-{Stealth[length=2.2mm]}, thick, orange!85!black]
  (lar) -- node[lbl, midway]{CO$_2$, calor,\\ vapor} (air);
\draw[{Stealth[length=2.2mm]}-{Stealth[length=2.2mm]}, thick, orange!85!black, dashed]
  (sus.south) .. controls ([xshift=1.4cm,yshift=-1.1cm]sus.south)
    and ([xshift=-2.0cm,yshift=-1.2cm]air.south) ..
  node[lbl, pos=0.5]{evaporación, CH$_4$} (air.south);

% Inlet (rótulo a la izquierda, flecha despejada)
\draw[flujo, green!50!black] ([yshift=1.3cm]box.west) ++(-2.0,0)
  -- node[pos=0, above=2pt, anchor=south west, font=\footnotesize, align=left]
     {inlet:\\ $T_{\mathrm{in}},\; w_{\mathrm{in}},\; Q$}
  ([yshift=1.3cm]box.west);

% Outlet (rótulo a la derecha, flecha despejada)
\draw[flujo, red!70!black] ([yshift=1.3cm]box.east)
  -- node[pos=1, above=2pt, anchor=south east, font=\footnotesize, align=right]
     {outlet: $Q$}
  ++(2.0,0);

% Sensores en el outlet (debajo de la flecha)
\node[sens] (sensores) at ([xshift=1.0cm,yshift=0.55cm]box.east) {$\mathbf{y}$};
\draw[-{Stealth[length=2mm]}, thick, red!70!black, dashed]
  (sensores.north) -- ([xshift=1.0cm]box.east);
\node[below=3pt of sensores, align=center, font=\footnotesize]
  {$T,\ H$\\ $\mathrm{CO_2},\ \mathrm{CH_4}$};

\end{tikzpicture}
\caption{El sistema completo de bioconversión instrumentado
(Definición~\ref{def32}). Dentro del volumen de control interactúan tres
grupos de compartimentos: sustrato, población larval y aire interno. La
frontera es atravesada por los flujos de ventilación (\emph{inlet/outlet},caudal $Q$); las mediciones $\mathbf{y}$ se toman en
la corriente de salida y una celda de carga registra la masa total del lote.}
\label{fig:sistema-completo}
\end{figure}

\subsection{Condiciones iniciales y horizonte del ciclo}

\begin{definition}[Ciclo de bioconversión]\label{def33}
Un \emph{ciclo de bioconversión} es la trayectoria del sistema
\eqref{eq:estado-completo} en el horizonte $[0, t_f]$, con condiciones
iniciales
\begin{equation}
N(0) = N_0, \quad S(0) = S_0, \quad
A(0) = A_0, \quad B(0) = B_0, \quad L(0) = L_0,
\label{eq:cond-iniciales}
\end{equation}
y ambiente en equilibrio con el \emph{inlet},
$T(0) = T_{\mathrm{in}}$, $w(0) = w_{\mathrm{in}}$,
$c_i(0) = c_{i,\mathrm{in}}$. El ciclo termina en $t_f$, definido como el
instante en que la fracción de prepupas en la población supera un umbral
$\pi_f$ convenido (típicamente $\pi_f \approx 0.5$), o equivalentemente
cuando la tasa agregada de asimilación cae por debajo del umbral de
mantenimiento de la población,
\begin{equation}
t_f \;=\; \inf\bigl\{ t > 0 \;:\; R_A(t) \;\leq\; N(t)\,m\,B(t) \bigr\}.
\label{eq:tf}
\end{equation}
\end{definition}

\begin{remark}[Criterios de terminación y su medición]\label{rem:tf}
Los dos criterios de \eqref{eq:tf} no son equivalentes operativamente: la
fracción de prepupas requiere inspección visual o muestreo, mientras que el
criterio metabólico es \emph{observable en línea} a través de la señal de
CO$_2$ ---cuando la producción total se reduce a la componente de
mantenimiento, $r_{\mathrm{CO_2}} \approx r_{\mathrm{CO_2},m}$, el
crecimiento y la acumulación de lípidos han cesado. Esta segunda forma es la
que el gemelo digital puede evaluar en tiempo real sin intervenir el lote.
Para larvas neonatas de cría sincronizada, los valores iniciales
$A_0, B_0, L_0$ corresponden al peso seco promedio de una larva de un día,
repartido entre compartimentos según su composición bioquímica medida.
\end{remark}